
% GA-PRAC-01
% C1 - Caracterizacion de los modelos de proceso
% Responsable: Cordova Carreno Mayra Lucila

\section{C1. Caracterización de los modelos de proceso de requisitos}

\subsection{Modelo secuencial o cascada}

El modelo secuencial o cascada organiza el desarrollo del software mediante fases principalmente consecutivas. Las actividades de Ingeniería de Requisitos se realizan de manera planificada y sus resultados sirven como base para las etapas posteriores.

\textbf{Fases del desarrollo:} análisis de requisitos, diseño, implementación, pruebas, despliegue y mantenimiento. Las actividades de elicitación, análisis, especificación y validación permiten establecer los requisitos que orientan el desarrollo. La gestión de requisitos acompaña el proceso.

\textbf{Supuestos:} se considera que los requisitos pueden identificarse y documentarse con suficiente estabilidad durante las etapas iniciales. Los cambios posteriores deben gestionarse mediante procedimientos de control.

\textbf{Condiciones de aplicación:} resulta favorable cuando los requisitos son relativamente estables, existen obligaciones documentales, se necesitan aprobaciones formales y las actividades deben seguir una planificación definida.

\textbf{Limitaciones:} los cambios tardíos pueden provocar retrabajo y afectar los costos y plazos. Además, una validación funcional realizada demasiado tarde puede dificultar la detección oportuna de errores.

\subsection{Modelo iterativo e incremental}

El modelo iterativo e incremental organiza el desarrollo mediante ciclos sucesivos. En cada iteración se revisan y perfeccionan los requisitos, mientras que los incrementos permiten incorporar progresivamente funcionalidades al sistema.

\textbf{Fases o actividades:} planificación de la iteración, elicitación y análisis de requisitos, especificación, diseño, implementación, pruebas y evaluación de resultados. Estas actividades pueden repetirse durante diferentes ciclos.

\textbf{Supuestos:} los requisitos pueden comprenderse y perfeccionarse progresivamente mediante la evaluación de resultados parciales y la retroalimentación de las partes interesadas.

\textbf{Condiciones de aplicación:} resulta favorable cuando existe incertidumbre sobre algunos requisitos, se necesitan validaciones progresivas y es posible organizar el sistema en incrementos funcionales.

\textbf{Limitaciones:} requiere coordinación entre iteraciones, gestión de dependencias, control de cambios y planificación de las funcionalidades que se incorporarán en cada incremento.

\subsection{Modelo ágil}

El modelo ágil se caracteriza por su capacidad de adaptación, la colaboración frecuente con las partes interesadas y la entrega continua o frecuente de valor. Los requisitos se refinan y priorizan conforme evoluciona el conocimiento del problema.

\textbf{Fases o actividades:} identificación de necesidades, elicitación continua, priorización y refinamiento de requisitos, planificación de ciclos cortos, desarrollo, pruebas y validación frecuente. Estas actividades no necesariamente siguen una secuencia rígida.

\textbf{Supuestos:} los requisitos pueden cambiar durante el proyecto y existe disponibilidad suficiente de las partes interesadas para proporcionar retroalimentación y participar en las decisiones.

\textbf{Condiciones de aplicación:} resulta favorable cuando existe alta volatilidad de requisitos, se necesita responder rápidamente a cambios y es posible mantener una colaboración frecuente con los stakeholders.

\textbf{Limitaciones:} la participación insuficiente de los stakeholders, la falta de priorización y un control inadecuado del alcance pueden dificultar el cumplimiento de los objetivos. Los proyectos con exigencias contractuales o normativas requieren mecanismos adicionales de documentación y trazabilidad.

\subsection{Comparación general}

El modelo secuencial prioriza la planificación y la definición temprana de requisitos; el modelo iterativo e incremental permite perfeccionarlos y validar resultados progresivamente; mientras que el modelo ágil enfatiza la adaptación continua y la colaboración frecuente.

La elección del modelo debe considerar la estabilidad de los requisitos, la disponibilidad de las partes interesadas, la criticidad del sistema, los riesgos y las obligaciones documentales del proyecto.
